\section*{अध्याय \ref{ch:b1:cross-sections} --- अनुप्रस्थ काट}

\begin{solution}{exo:b1:cross-sections:1}
1:25\,000 पर मानचित्र का \qty{1}{cm} ज़मीन पर \qty{250}{m} है। ऊर्ध्वाधर पैमाना \qty{100}{m} के लिए \qty{1}{cm} है, अतः
ऊर्ध्वाधर पैमाना क्षैतिज का $250/100 = 2.5$ गुना है: $e = 2.5$।
\end{solution}

\begin{solution}{exo:b1:cross-sections:2}
$\gamma = \ang{90} - \ang{60} = \ang{30}$, और $\tan\beta = \tan\ang{40}
\cos\ang{30} = 0.8391\times0.8660 = 0.7267$: $\beta = \ang{36.0}$, परत काट के पूर्वी सिरे की ओर नीचे जाती है। दिगंश
\ang{000} की काट में $\gamma = \ang{90}$ और $\beta = 0$ है: काट \omterm{def:b1:geological-maps:attitude}{स्ट्राइक} के अनुदिश चलती है और परत क्षैतिज दिखती है।
\end{solution}

\begin{solution}{exo:b1:cross-sections:3}
$t = h\cos\delta = 120\times\cos\ang{12.16} = 120\times0.9776 =
\qty{117.3}{m}$।
\end{solution}

\begin{solution}{exo:b1:cross-sections:4}
$\tan\ang{12} = 0.2126$। $e = 2$ के साथ: $\arctan 0.4251 = \ang{23.0}$; $e = 3$ के साथ:
$\arctan 0.6377 = \ang{32.5}$; $e = 5$ के साथ: $\arctan 1.0628 =
\ang{46.7}$। कोण $e$ से कहीं कम बढ़ता है, और कभी \ang{90} तक नहीं पहुँचता।
\end{solution}

\begin{solution}{exo:b1:cross-sections:5}
(क) $t = 300\sin\ang{25} = \qty{126.8}{m}$। (ख) $t = 300\sin\ang{35} =
\qty{172.1}{m}$। (ग) $t = 300\sin\ang{15} = \qty{77.6}{m}$। एक ही \omterm{def:b1:geological-maps:geological-map}{बहिर्भाग} चौड़ाई का अर्थ दो गुने से अधिक भिन्न
मोटाइयाँ हो सकता है: ज़मीन की ढलान को सदा ध्यान में रखना होगा।
\end{solution}

\begin{solution}{exo:b1:cross-sections:6}
\omterm{def:b1:geological-maps:attitude}{नति-दिशा} $\alpha$ और \omterm{def:b1:geological-maps:attitude}{नति} $\delta$ के साथ, प्रमेय उत्तर-दक्षिण काट ($\theta = 0$) में $\tan\ang{20} = \tan\delta\cos\alpha$
देता है, और पूर्व-पश्चिम काट ($\theta = \ang{90}$) में $\tan\ang{30} =
\tan\delta\cos(\alpha - \ang{90}) = \tan\delta\sin\alpha$। भाग देने पर
$\tan\alpha = 0.5774/0.3640 = 1.586$, $\alpha = \ang{057.8}$; वर्गों को जोड़ने पर,
$\tan\delta = \sqrt{0.3640^2 + 0.5774^2} = 0.6825$, $\delta =
\ang{34.3}$। दोनों आभासी नतियाँ वास्तविक \omterm{def:b1:geological-maps:attitude}{नति} से छोटी हैं, जैसा होना ही चाहिए।
\end{solution}

\begin{solution}{exo:b1:cross-sections:7}
$100/\tan\ang{12.16} = 100/0.2155 = \qty{464.1}{m}$। मानचित्र पर पूर्वी भुजा के \omterm{def:b1:cross-sections:structure-contour}{संरचना समोच्च} सचमुच आधे किलोमीटर
से थोड़े कम दूर हैं, समान अंतराल पर और समांतर: भुजा एक समतल है।
\end{solution}

\begin{solution}{exo:b1:cross-sections:8}
B1: $524.02 - 488.98 = \qty{35.04}{m}$; B3: $529.98 - 494.93 =
\qty{35.05}{m}$। B2 के ऊपर शेल का आधार चूना-पत्थर का ऊपरी सिरा है, जो उसके आधार से \qty{180}{m} ऊपर है:
$280.16 + 180 = \qty{460.16}{m}$, जो ज़मीन (\qty{380.64}{m}) से ऊपर है, इसीलिए वहाँ वह अपरदित है।
B1 और B2 के बीच \omterm{def:b1:geological-maps:geological-map}{संपर्क} \qty{2000}{m} पर $460.16 - 35.04 = \qty{425.12}{m}$ ऊपर चढ़ता है: $0.2126$ की ढलान,
\ang{12.0} की \omterm{def:b1:cross-sections:apparent-dip}{आभासी नति}। B1 और B3, B2 के दोनों ओर समान दूरी पर, हर \omterm{def:b1:geological-maps:geological-map}{संपर्क} को \qty{1}{cm} के भीतर एक ही उन्नयन पर
पाते हैं: \omterm{def:b1:geological-maps:fold}{वलन} $x = \qty{3.5}{km}$ पर एक ऊर्ध्वाधर अक्षीय तल के सापेक्ष सममित है, उसकी दोनों भुजाएँ समान रूप से
झुकती हैं, और उनके बीच कोई भ्रंश काट को नहीं काटता।
\end{solution}

\begin{solution}{exo:b1:cross-sections:9}
\omterm{def:b1:cross-sections:throw}{विस्थापन} $= 150/\tan\ang{50} = 150/1.1918 = \qty{125.9}{m}$। \ang{150} दिगंश की काट में $\gamma = \ang{30}$ और
$\tan\beta = 1.1918\times0.8660 =
1.0321$: भ्रंश \ang{45.9} की \omterm{def:b1:cross-sections:apparent-dip}{आभासी नति} के साथ दिखता है। क्षैतिज परत के दो कटाव-सिरे उस रेखा पर ऊर्ध्वाधर रूप से
\qty{150}{m} अलग हैं, अतः काट के अनुदिश $150/1.0321 = \qty{145.3}{m}$ अलग हैं, जो \omterm{def:b1:cross-sections:throw}{क्षैतिज विस्थापन} से अधिक है,
क्योंकि काट भ्रंश को तिरछा पार करती है।
\end{solution}

\begin{solution}{exo:b1:cross-sections:10}
$\cos\gamma = \tan\ang{20}/\tan\ang{40} = 0.3640/0.8391 = 0.4338$, इसलिए
$\gamma = \ang{64.3}$: दिगंशों $\alpha \pm \ang{64.3}$ की काटें (और विपरीत दिशाएँ, जहाँ परत ऊपर उठती है)। \omterm{def:b1:cross-sections:apparent-dip}{आभासी नति}
$\gamma = 0$, \ang{10}, \ang{20}, \ang{30} के लिए \ang{40.0}, \ang{39.6}, \ang{38.3} और \ang{36.0} है, परंतु
\ang{22.8}, \ang{16.0}, \ang{8.3} और $0$ क्रमशः $\gamma = \ang{60}$, \ang{70}, \ang{80}, \ang{90} के लिए। सचमुच, $\tan\beta = \tan\delta\cos\gamma$ का अवकलन करने पर,
$\dfrac{\mathrm d\beta}{\mathrm d\gamma} =
-\dfrac{\tan\delta\sin\gamma}{1 + \tan^2\delta\cos^2\gamma}$, जो $\gamma = 0$ पर शून्य है और
$-\tan\delta$ के बराबर होता है जब $\gamma = \ang{90}$: \omterm{def:b1:geological-maps:attitude}{नति-दिशा} से थोड़ी हटी काट लगभग वास्तविक \omterm{def:b1:geological-maps:attitude}{नति} दिखाती है, जबकि \omterm{def:b1:geological-maps:attitude}{स्ट्राइक} के निकट दिगंश का
हर अंश \omterm{def:b1:cross-sections:apparent-dip}{आभासी नति} को लगभग एक अंश बदल देता है।
\end{solution}

\begin{solution}{exo:b1:cross-sections:11}
परतें पूर्व की ओर झुकती हैं, अतः परत-तल का अभिलंब, नीचे की ओर अनुसरित, ऊर्ध्वाधर से \ang{12.16} पश्चिम की ओर,
\ang{260.7} की ओर झुका है। (क) एक ऊर्ध्वाधर बोरहोल उसके साथ $\psi = \ang{12.16}$ कोण बनाता है:
$176.0/\cos\ang{12.16} = \qty{180.0}{m}$, \omterm{def:b1:cross-sections:thickness}{ऊर्ध्वाधर मोटाई}। (ख) \ang{260.7} की ओर \ang{60} नमन करता बोरहोल, अभिलंब
की उसी ओर, ऊर्ध्वाधर से \ang{30} झुका है: $\psi = \ang{30} - \ang{12.16} =
\ang{17.84}$, $\cos\psi = 0.9519$, और कोर $176.0/0.9519 =
\qty{184.9}{m}$ लंबा है। (ग) \ang{080.7} की ओर यह दूसरी ओर झुका है: $\psi = \ang{30} + \ang{12.16} = \ang{42.16}$,
$\cos\psi = 0.7413$, और कोर $176.0/0.7413 = \qty{237.4}{m}$ लंबा है। कोर सबसे छोटा, $t$ के बराबर, तब होता है जब
बोरहोल अभिलंब के अनुदिश चलता है, यानी \ang{260.7} की ओर \ang{77.84} नमन करता हुआ।
\end{solution}

\begin{solution}{exo:b1:cross-sections:12}
पूर्व-पश्चिम ऊर्ध्वाधर समतल में, कोण पूर्व से मापे गए हैं: ऊपरी परतें पूर्व की ओर \ang{20} उठती हैं और निचली परतें
पूर्व की ओर \ang{70} नीचे जाती हैं। \omterm{def:b1:geological-maps:geological-map}{बहिर्भाग} को \omterm{def:b1:geological-maps:attitude}{स्ट्राइक} के परितः \ang{20} घुमाने पर, पूर्वी ओर नीचे, ऊपरी परतें क्षैतिज
हो जाती हैं और निचली परतें $\ang{70} + \ang{20} = \ang{90}$ पर आ जाती हैं: जब ऊपरी परतें जमा हुईं, तब वे ऊर्ध्वाधर
थीं। इतिहास, सबसे पुरानी घटना से: निचली परतें समतल जमा हुईं; वे किनारे पर, ऊर्ध्वाधर तक, मोड़ी गईं और उठाई गईं; उनके
किनारे एक सतह तक अपरदित हुए; ऊपरी परतें उस सतह के आरपार जमा हुईं; पूरा पश्चिम की ओर \ang{20} नीचे झुक गया; अपरदन ने
आज की ज़मीन की सतह काटी। यह हटन के सिकार पॉइंट के \omterm{def:b1:geological-maps:geological-map}{बहिर्भाग} का इतिहास है।
\end{solution}

\begin{solution}{pb:b1:cross-sections:1}
\textbf{1.} कर से ऊपर की चट्टान की मोटाई $z - 200$ है: $x = \qty{1}{km}$ पर सबसे बड़ी, $556 - 200 = \qty{356}{m}$;
$x =
\qty{9}{km}$ पर सबसे छोटी, $339 - 200 = \qty{139}{m}$।

\textbf{2.} \qty{1}{km} उन्नयन \qty{4.2}{cm} मापता है और \qty{1}{km} दूरी \qty{1.4}{cm}: $e = 4.2/1.4 = 3$।

\textbf{3.} $380.64 - 220.48 = \qty{160.16}{m}$।

\textbf{4.} B2 में चूना-पत्थर का आधार $380.64 - 100.48 =
\qty{280.16}{m}$ पर है और बलुआ पत्थर का आधार \qty{160.16}{m} पर: \qty{200}{m} पर सुरंग बलुआ पत्थर में है।

\textbf{5.} B1: $524.02 - 668.98 = \qty{-144.96}{m}$; B3: $529.98 - 674.93 =
\qty{-144.95}{m}$।

\textbf{6.} पश्चिम से पूर्व: खड़िया, शेल, चूना-पत्थर, बलुआ पत्थर (\omterm{def:b1:geological-maps:fold}{अपनति} के क्रोड में), चूना-पत्थर, शेल, और फिर खड़िया
$x = \qty{9}{km}$ तक, \omterm{def:b1:geological-maps:fold}{अभिनति} से होकर और भ्रंश के आरपार।

\textbf{7.} $\gamma = \ang{80.7} - \ang{90} = \ang{-9.3}$; $\tan\beta =
\tan\ang{12.16}\cos\ang{9.3} = 0.2155\times0.9869 = 0.2126$: $\beta =
\ang{12.0}$, पूर्व की ओर नीचे।

\textbf{8.} $(280.16 + 144.95)/2000 = 425.11/2000 = 0.2126$: \ang{12.0}।

\textbf{9.} $\arctan(3\times0.2126) = \arctan 0.6377 = \ang{32.5}$।

\textbf{10.} $\gamma = \ang{120} - \ang{90} = \ang{30}$; $\tan\beta =
\tan\ang{60}\cos\ang{30} = 1.732\times0.866 = 1.500$: $\beta = \ang{56.3}$।

\textbf{11.} भुजाएँ \ang{080.7} या \ang{279.3} की ओर झुकती हैं, काट की रेखा से केवल \ang{9.3} दूर, और $\cos\ang{9.3} = 0.987$;
भ्रंश \ang{120} की ओर झुकता है, उससे \ang{30} दूर, और $\cos\ang{30} = 0.866$।

\textbf{12.} भ्रंश-तल $z = 250 - \tan\ang{60}\,[(x - 4300)
\sin\ang{120} + (y - 1000)\cos\ang{120}]$ है। $y = 6000$ के साथ $(y - 1000)
\cos\ang{120} = -2500$, अतः $(x - 4300)\sin\ang{120} = 2500 - (z - 250)/
\tan\ang{60}$, और $1/\tan\ang{60} = \tan\ang{30}$।

\textbf{13.} $z = 200$ पर: $x = 4300 + (2500 + 50\times0.5774)/0.8660 =
4300 + 2528.9/0.8660 = \qty{7220}{m}$, यानी C से \qty{7.22}{km}।

\textbf{14.} \omterm{def:b1:cross-sections:throw}{क्षैतिज विस्थापन} $= 200/\tan\ang{60} = \qty{115.5}{m}$; C--D के अनुदिश अलगाव $= 200/1.500 = \qty{133.3}{m}$।

\textbf{15.} भ्रंश पूर्व-दक्षिण-पूर्व की ओर झुकता है और सामान्य है: लटकती दीवार, पूर्व में, नीचे खिसकी है। सुरंग के स्तर पर
खड़िया भ्रंश के दोनों ओर है: पश्चिम में \omterm{def:b1:geological-maps:fold}{अभिनति} की खड़िया; पूर्व में वह खड़िया जो \qty{200}{m} नीचे लाई गई है।

\textbf{16.} सुरंग भ्रंश के आरपार खड़िया में रहती है, इसलिए मुख पर चट्टान का कोई परिवर्तन उसकी घोषणा नहीं करता; केवल
पिसी चट्टान का एक क्षेत्र, और शायद पानी, करता है। इंजीनियर उसकी स्थिति काट से अनुमानित करते हैं (यहाँ
\qty{7.22}{km}), मशीन के पास आने पर उसके आगे जाँच-छेद खोदते हैं, और पिसी चट्टान के क्षेत्र के लिए सहारा तैयार करते हैं।

\textbf{17.} $668.98 - 488.98 = \qty{180.00}{m}$।

\textbf{18.} $t = 180\cos\ang{12.16} = \qty{176.0}{m}$।

\textbf{19.} कटक ($x = 3500$) पर बलुआ पत्थर का आधार \qty{160.16}{m} पर है (प्रश्न 3); वह दोनों भुजाओं पर प्रश्न 7 की
\ang{12.0} की \omterm{def:b1:cross-sections:apparent-dip}{आभासी नति} के साथ नीचे जाता है, यानी अक्ष से क्षैतिज दूरी के प्रति मीटर $\tan\ang{12.0}$ से।

\textbf{20.} चूना-पत्थर का आधार बलुआ पत्थर के आधार से \qty{120}{m} ऊपर है, $280.16 - u\tan\ang{12.0}$ पर, जहाँ $u = |x - 3500|$,
और उसका ऊपरी सिरा \qty{180}{m} और ऊपर, $460.16 - u\tan\ang{12.0}$ पर। $\tan\ang{12.0} =
0.21256$ के साथ: आधार \qty{200}{m} से नीचे है जब $u > 80.16/0.21256 = 377.1$, और ऊपरी सिरा उससे ऊपर है जब
$u < 260.16/0.21256 = 1224.0$।

\textbf{21.} $x = 3500 \pm 377.1$ और $3500 \pm 1224.0$: सुरंग \qty{2276}{m} से \qty{3123}{m} तक और \qty{3877}{m} से
\qty{4724}{m} तक चूना-पत्थर में है।

\textbf{22.} एक क्षैतिज रेखा के अनुदिश, चूना-पत्थर का आधार और ऊपरी सिरा, ऊर्ध्वाधर रूप से \qty{180}{m} अलग और दोनों
$\tan\beta$ पर ढलते हुए, क्षैतिज दूरी $h/\tan\beta$ पर पार किए जाते हैं: $180/0.21256 = \qty{846.8}{m}$ हर भुजा पर,
पूर्णांकन तक $1224.0 - 377.1 = 846.9$ के बराबर।

\textbf{23.} $2\times846.8 = \qty{1693.7}{m}$: सुरंग \textbf{लगभग \qty[text-series-to-math=true]{1694}{m} चूना-पत्थर} खोदती है,
जो उसकी \qty{8}{km} लंबाई के पाँचवें भाग से अधिक है, \omterm{def:b1:geological-maps:fold}{अपनति} के बलुआ पत्थर क्रोड के दोनों ओर \qty{847}{m} के
दो हिस्सों में।
\end{solution}
